% 01-vi-sao-chuoi-khac.tex - Vì sao dữ liệu chuỗi khác dữ liệu thường
%
% Scope: what makes a sequence different from i.i.d. data; examples from
% text, speech, video, time series; the memory requirement.

\section{Vì sao dữ liệu chuỗi khác}
\label{sec:ch01-vi-sao-chuoi-khac}

Một bức ảnh chụp một con mèo. Nếu bạn xáo trộn các pixel của nó, nó vẫn là
một con mèo; mô hình phân loại ảnh không cần biết pixel nào đứng cạnh pixel
nào, chỉ cần biết tất cả chúng cùng có mặt. Đó là thế giới của dữ liệu
\tn{độc lập và phân phối giống hệt nhau}{independent and identically
distributed} (i.i.d.), nơi mỗi mẫu huấn luyện là một điểm rời rạc không liên
quan đến mẫu trước nó.

Giờ thử xáo trộn các chữ trong câu bạn vừa đọc.

Chuỗi thì khác. Trật tự mang thông tin, và thông tin đó không nằm trong từng
phần tử riêng lẻ mà nằm ở quan hệ giữa chúng. Một mô hình đọc ``mèo đuổi
chuột'' phải biết rằng ``mèo'' đứng trước ``đuổi'' và ``đuổi'' đứng trước
``chuột''; đảo thành ``chuột đuổi mèo'' thì nghĩa đảo ngược hoàn toàn, dù
tập hợp từ không đổi.

Dữ liệu dạng \tn{chuỗi}{sequence} có mặt ở khắp nơi. Văn bản là chuỗi các từ
hoặc ký tự. Tiếng nói là chuỗi các mẫu âm thanh lấy ở tần số lấy mẫu cố định.
Video là chuỗi các khung hình. Dữ liệu tài chính là chuỗi giá theo thời gian.
DNA là chuỗi các nucleotide. Trong mọi trường hợp, ý nghĩa của một phần tử
phụ thuộc vào những phần tử đứng trước nó; và đôi khi cả những phần tử đứng
sau.

Để xử lý dữ liệu chuỗi, một mô hình cần \textbf{trí nhớ}. Nó phải giữ lại
thông tin từ quá khứ để diễn giải hiện tại. Một mạng nơ-ron truyền thẳng
thông thường không làm được việc này: mỗi đầu vào được xử lý độc lập, đầu ra
chỉ phụ thuộc vào đầu vào hiện tại. Cần một kiến trúc cho phép tín hiệu đi
từ bước thời gian trước sang bước thời gian sau.

Đó chính là \tn{mạng hồi quy}{recurrent neural network} (RNN). Kiến trúc này
giữ một \tn{trạng thái ẩn}{hidden state}; một vector tổng kết mọi thứ mạng
đã thấy cho đến thời điểm hiện tại. Tại mỗi bước, trạng thái ẩn được cập nhật
bằng cách kết hợp đầu vào mới với trạng thái ẩn cũ. Cùng một tập trọng số
được dùng lại ở mọi bước thời gian, nên mạng có thể xử lý chuỗi có độ dài bất
kỳ.

Nhưng trí nhớ có giá của nó. Mạng hồi quy không học bằng cách truyền thẳng
rồi lan truyền ngược như mạng nơ-ron thường. Vì trạng thái ẩn tại bước $t$
phụ thuộc vào trạng thái ẩn tại bước $t-1$, gradient của hàm mất mát phải
được truy ngược qua \textbf{từng bước thời gian}, từ cuối chuỗi về đầu chuỗi.
Kỹ thuật này gọi là \tn{lan truyền ngược qua thời gian}{backpropagation
through time} (BPTT), và nó là trung tâm của mọi thứ chương này sẽ làm.

Mục tiêu của chương này là ba thứ. Thứ nhất: thiết lập hệ ký hiệu mà toàn bộ
cuốn sách sẽ dùng. Thứ hai: dẫn xuất BPTT đến từng Jacobian một, để ba chương
tiếp theo có thể mổ xẻ chính xác vì sao gradient biến mất hoặc bùng nổ. Thứ
ba: chạy một ví dụ bằng số để bạn thấy mọi ma trận, mọi đạo hàm, và để bạn
kiểm tra lại gradient bằng \tn{sai phân hữu hạn}{finite difference}; cách
duy nhất để biết chắc một cài đặt lan truyền ngược là đúng.
